\begin{center}
    {\Large\itshape Alban pendant la guerre de Cent Ans}
\end{center}

\vspace{1.5em}

% Le bord gauche du scan est rogné : les débuts de plusieurs lignes
% ont été reconstitués lorsque le contexte et les fragments visibles le permettent.
\noindent Peu de renseignements nous sont parvenus sur la 1\textsuperscript{ère} partie du XIV\textsuperscript{ème} siècle. On note au passage qu'en 1342, un « Frédalen de Lumbenchir »\footnote{Aucune trace de cette personne n'est retrouvée dans les manuels classiques d'histoire. Est-ce une mauvaise transcription, ou bien un nom local que l'on ne trouve que dans les archives locales? Cela mérite d'être tiré au clair.} est seigneur d'Alban. Est-ce par ses soins qu'Alban s'organise en poste avancé en sentinelle de la région albigeoise ? La chose est possible. En fait, on n'a pas seulement à se défendre contre les Anglais mais aussi contre des bandes de maquisards multiples. Les comtes d'Armagnac, maîtres du Rouergue, les comtes de Foix qui sont en même temps vicomtes de Lautrec, les ducs de Berry et d'Anjou, qui sont gouverneurs du Languedoc, vident leurs querelles au détriment des populations. Les troupes de ces grands seigneurs féodaux sont, avec quelques militaires de métier, un ramassis de pillards et de brigands. Ces armées hétéroclites occupent les châteaux qui leur serviront de repaires. Ils pillent les campagnes, détroussant les uns, malmenant les autres. C'est alors que le rôle d'Alban comme sentinelle avancée et gardienne de la paix se manifeste. Il faut organiser la défense derrière les remparts, le cas échéant, traiter avec le chef ennemi, obtenir son éloignement, ne fût-il que momentané pour faire la moisson ou les vendanges.

\vspace{0.8em}

\noindent En 1360, les consuls d'Albi envoient un homme de confiance aux consuls de Trébas pour s'informer du passage des 500 hommes d'armes du Comte d'Armagnac.
